\chapter*{Заключение} \label{ch-conclusion}
\addcontentsline{toc}{chapter}{Заключение}	% в оглавление 

\noindent
В данной выпускной квалификационной работе была разработана, программно 
реализована и протестирована модель глубокого обучения для прогнозирования 
фенотипических признаков сельскохозяйственных растений по генотипическим 
данным, представленным в виде однонуклеотидных полиморфизмов (SNP). 
Исследование проводилось на трёх культурах: нут (\textit{Cicer arietinum}), 
рис (\textit{Oryza sativa}) и рапс (\textit{Brassica napus}), что позволило 
оценить обобщающую способность предложенной архитектуры.

\subsection*{Основные результаты работы}

\begin{enumerate}
	\item \textbf{Разработана архитектура модели}, включающая четыре 
	основных блока: генетического кодирования с механизмом внимания 
	на уровне SNP, трансформерный блок для моделирования межгенных 
	взаимодействий, блок прогнозирования фенотипов и реконструкционный 
	блок для регуляризации скрытых представлений. Предложенная 
	архитектура позволяет адаптивно взвешивать вклад отдельных 
	SNP-маркеров при формировании генных представлений и учитывать 
	сложные нелинейные зависимости между генами.
	
	\item \textbf{Реализован программный комплекс} на базе библиотеки 
	PyTorch, включающий модули для загрузки и предобработки генетических 
	данных, частотного кодирования SNP, аугментации методом MixUp, 
	обучения модели с многозадачной функцией потерь и адаптивным 
	взвешиванием признаков, а также для анализа межгенных взаимодействий 
	с использованием значений Шепли.
	
	\item \textbf{Проведена серия вычислительных экспериментов} 
	на данных трёх культур. Наибольшая точность прогнозирования 
	достигнута для времени цветения рапса (коэффициент корреляции 
	Пирсона 0.809) и массы тысячи семян нута (0.881). Для данных 
	риса получены корреляции 0.830 для высоты растения и 0.516 
	для числа побегов.
	
	\item \textbf{Исследованы методы аугментации данных}. Сравнительный 
	анализ показал, что метод MixUp эффективнее аугментации перестановкой 
	хромосом, особенно для сложных метаболических признаков (содержание 
	глюкозинолатов: MAPE улучшился на 8.93\%), при этом для хорошо 
	предсказуемых признаков (время цветения) аугментация не требуется.
	
	\item \textbf{Сравнены стратегии отбора генов} для работы с 
	высокоразмерными геномными данными. Равномерный отбор по геному 
	и равномерный отбор внутри хромосом существенно превосходят 
	взвешенный отбор по качеству предсказания времени цветения 
	(0.850 против 0.798), при этом требуя на 26\% меньше SNP.
	
	\item \textbf{Выполнена интерпретация модели} с помощью значений 
	Шепли. На данных риса выявлены гены-«хабы» \texttt{Os01g0125800} 
	и \texttt{Os01g0150000}, участвующие в наибольшем количестве 
	взаимодействий. На данных рапса ключевым геном-«хабом» является 
	\texttt{gene-LOC106381653}. Структура выявленных взаимодействий 
	(наличие генов с множеством как положительных, так и отрицательных 
	связей) соответствует паттерну ключевых регуляторных генов.
\end{enumerate}

\subsection*{Научная новизна и практическая значимость}

\textbf{Научная новизна} работы заключается в предложении архитектуры 
модели с механизмом внимания на уровне SNP и реконструкционной 
регуляризацией, а также в разработке многозадачной функции потерь 
с адаптивным взвешиванием фенотипических признаков.

\textbf{Практическая значимость} состоит в создании программной 
реализации модели, которая может быть использована в селекционных 
программах для идентификации генетических маркеров \cite{Atieno2017, Bankin2024}, связанных с 
продуктивностью, устойчивостью к засухе и другими хозяйственно 
ценными признаками. Полученные результаты подтверждают применимость 
подхода в регионах рискованного земледелия для ускоренного выведения 
адаптированных сортов.

\subsection*{Перспективы дальнейших исследований}

Дальнейшее развитие работы может включать следующие направления:

\begin{enumerate}
	\item Тестирование модели на других сельскохозяйственных культурах 
	(пшеница, ячмень, соя) для подтверждения обобщающей способности.
	
	\item Экспериментальная валидация выявленных генов-«хабов» 
	(\texttt{Os01g0125800} у риса и \texttt{gene-LOC106381653} у рапса) 
	с использованием методов обратной генетики или анализа экспрессии.
	
	\item Разработка веб-сервиса для онлайн-прогнозирования фенотипов 
	по SNP-данным, доступного для селекционеров и исследователей.
	

\end{enumerate}

\subsection*{Итоговый вывод}

Разработанная модель глубокого обучения на основе трансформерной 
архитектуры с механизмом внимания на уровне SNP и реконструкционной 
регуляризацией успешно справляется с задачей прогнозирования 
фенотипических признаков сельскохозяйственных культур. Эффективность 
подхода подтверждена на данных трёх культур, различающихся по 
плотности SNP-маркеров и геномной архитектуре. Возможность 
интерпретации модели с помощью значений Шепли позволяет выявлять 
значимые межгенные взаимодействия и гены-кандидаты, что делает 
разработанный инструмент ценным не только для прогнозирования, 
но и для понимания молекулярных механизмов формирования 
агрономически важных признаков.

%Заключение (2 -- 5 страниц) обязательно содержит выводы по теме работы, \textit{конкретные
%предложения и рекомендации} по исследуемым вопросам. Количество общих выводов
%должно вытекать из количества задач, сформулированных во введении выпускной
%квалификационной работы.

%Предложения и рекомендации должны быть органически увязаны с выводами
%и направлены на улучшение функционирования исследуемого объекта. При разработке
%предложений и рекомендаций обращается внимание на их обоснованность,
%реальность и практическую приемлемость.

%Заключение не должно содержать новой информации, положений, выводов и
%т. д., которые до этого не рассматривались в выпускной квалификационной работе.
%Рекомендуется писать заключение в виде тезисов.

%Последним абзацем в заключении можно выразить благодарность всем людям, которые помогали автору в написании ВКР.